% GENERATED FILE. Edit canonical lessons, then run npm run generate:books.
% canonical-source-hash: fnv1a64:9df504382025247b
% canonical-lessons: LA-C232-furor, LA-C232-detego, LA-C232-videor, LA-C232-servio, LA-C232-patior

\chapter{To steal, To uncover, To seem, To serve, To suffer}
\label{ch:la-to-steal-to-uncover-to-seem-to-serve-to-suffer}
\hlchaptermodality{232}

\begin{chapteropening}
\textbf{By the end of this chapter:} \emph{I can say to steal, to uncover, to seem, to serve and to suffer.}

Complete the last lesson of chapter 232: I can say to steal, to uncover, to seem, to serve and to suffer.
\end{chapteropening}

\section[fūror, fūrārī]{fūror, fūrārī — to steal}

\label{lesson:LA-C232-furor}



\begin{quote}
Before the new one: say the Latin for to complain, then the Latin for to be missing.
\end{quote}

\subsection*{You'll want to know: fūror, fūrārī}
\textbf{fūror, fūrārī} — \textquotedblleft{}to steal\textquotedblright{}.

\begin{cousinweb}
\textbf{fūror} is built on \emph{fūr}, \textquotedblleft{}a thief\textquotedblright{} (Lewis \& Short), usually traced to Proto-Indo-European \emph{*bhōr}, \textquotedblleft{}thief\textquotedblright{} (EtymDB 2.1), from the root of \emph{ferō}, \textquotedblleft{}to carry\textquotedblright{}: one who carries off. It is a different word from \emph{furor}, \textquotedblleft{}rage\textquotedblright{}.
\end{cousinweb}

\begin{grammarlens}[title={What you've built}]
One more word for this chapter.
\end{grammarlens}

\subsection*{Your turn}
\begin{itemize}
\raggedright
  \item \emph{Say it:} \emph{fūror, fūrārī}
  \item \emph{Say it:} \emph{fūror, fūrārī}, once more
  \item \emph{Say it:} say \emph{fūror, fūrārī}
  \item \emph{Recall:} say the Latin for to complain, then the Latin for to be missing, then say \emph{fūror, fūrārī} again
\end{itemize}

\begin{tcolorbox}[breakable,colback=teal!4,colframe=teal!35!black,title={Before you move on}]
What does fūror, fūrārī mean? (\textquotedblleft{}To steal\textquotedblright{}.) Say it once more.
\end{tcolorbox}
\section[dētegō, dētegere]{dētegō, dētegere — to uncover, to discover}

\label{lesson:LA-C232-detego}



\begin{quote}
Before the new one: say the Latin for to be missing, then the Latin for to steal.
\end{quote}

\subsection*{You'll want to know: dētegō, dētegere}
\textbf{dētegō, dētegere} — \textquotedblleft{}to uncover, to discover\textquotedblright{}.

\begin{cousinweb}
\textbf{dētegō} is \emph{dē}, \textquotedblleft{}off\textquotedblright{}, with \emph{tegō}, \textquotedblleft{}to cover\textquotedblright{} (EtymDB 2.1): to take the cover off. English \emph{detect} comes from it.
\end{cousinweb}

\begin{grammarlens}[title={What you've built}]
One more word for this chapter.
\end{grammarlens}

\subsection*{Your turn}
\begin{itemize}
\raggedright
  \item \emph{Say it:} \emph{dētegō, dētegere}
  \item \emph{Say it:} \emph{dētegō, dētegere}, once more
  \item \emph{Say it:} say \emph{dētegō, dētegere}
  \item \emph{Recall:} say the Latin for to be missing, then the Latin for to steal, then say \emph{dētegō, dētegere} again
\end{itemize}

\begin{tcolorbox}[breakable,colback=teal!4,colframe=teal!35!black,title={Before you move on}]
What does dētegō, dētegere mean? (\textquotedblleft{}To uncover, to discover\textquotedblright{}.) Say it once more.
\end{tcolorbox}
\section[videor, vidērī]{videor, vidērī — to seem}

\label{lesson:LA-C232-videor}



\begin{quote}
Before the new one: say the Latin for to steal, then the Latin for to uncover.
\end{quote}

\subsection*{You'll want to know: videor, vidērī}
\textbf{videor, vidērī} — \textquotedblleft{}to seem\textquotedblright{}.

\begin{cousinweb}
\textbf{videor} is the passive of \emph{videō}, \textquotedblleft{}to see\textquotedblright{}: to be seen, so to seem. \emph{Videō} is inherited from Proto-Indo-European \emph{*weyd-}, \textquotedblleft{}to see, to know\textquotedblright{} (EtymDB 2.1).
\end{cousinweb}

\begin{grammarlens}[title={What you've built}]
One more word for this chapter.
\end{grammarlens}

\subsection*{Your turn}
\begin{itemize}
\raggedright
  \item \emph{Say it:} \emph{videor, vidērī}
  \item \emph{Say it:} \emph{videor, vidērī}, once more
  \item \emph{Say it:} say \emph{videor, vidērī}
  \item \emph{Recall:} say the Latin for to steal, then the Latin for to uncover, then say \emph{videor, vidērī} again
\end{itemize}

\begin{tcolorbox}[breakable,colback=teal!4,colframe=teal!35!black,title={Before you move on}]
What does videor, vidērī mean? (\textquotedblleft{}To seem\textquotedblright{}.) Say it once more.
\end{tcolorbox}
\section[serviō, servīre]{serviō, servīre — to serve, to be a slave}

\label{lesson:LA-C232-servio}



\begin{quote}
Before the new one: say the Latin for to uncover, then the Latin for to seem.
\end{quote}

\subsection*{You'll want to know: serviō, servīre}
\textbf{serviō, servīre} — \textquotedblleft{}to serve, to be a slave\textquotedblright{}.

\begin{cousinweb}
\textbf{serviō} is the verb made from \emph{servus}, \textquotedblleft{}a slave\textquotedblright{} (EtymDB 2.1): to be a slave, to serve. English \emph{serve} comes from it through French.
\end{cousinweb}

\begin{grammarlens}[title={What you've built}]
One more word for this chapter.
\end{grammarlens}

\subsection*{Your turn}
\begin{itemize}
\raggedright
  \item \emph{Say it:} \emph{serviō, servīre}
  \item \emph{Say it:} \emph{serviō, servīre}, once more
  \item \emph{Say it:} say \emph{serviō, servīre}
  \item \emph{Recall:} say the Latin for to uncover, then the Latin for to seem, then say \emph{serviō, servīre} again
\end{itemize}

\begin{tcolorbox}[breakable,colback=teal!4,colframe=teal!35!black,title={Before you move on}]
What does serviō, servīre mean? (\textquotedblleft{}To serve, to be a slave\textquotedblright{}.) Say it once more.
\end{tcolorbox}
\section[patior, patī]{patior, patī — to suffer, to allow}

\label{lesson:LA-C232-patior}



\begin{quote}
Before the new one: say the Latin for to seem, then the Latin for to serve.
\end{quote}

\subsection*{You'll want to know: patior, patī}
\textbf{patior, patī} — \textquotedblleft{}to suffer, to allow\textquotedblright{}.

\begin{cousinweb}
\textbf{patior} has no sure origin. Lewis \& Short compare Greek \emph{penthos}, \textquotedblleft{}grief\textquotedblright{}, and EtymDB 2.1 records a proposed root \emph{*pe(i)-}, \textquotedblleft{}to hurt\textquotedblright{}, but neither link is settled. English \emph{passion} and \emph{patient} come from it.
\end{cousinweb}

\begin{grammarlens}[title={What you've built}]
One more word for this chapter.
\end{grammarlens}

\subsection*{Your turn}
\begin{itemize}
\raggedright
  \item \emph{Say it:} \emph{patior, patī}
  \item \emph{Say it:} \emph{patior, patī}, once more
  \item \emph{Say it:} say \emph{patior, patī}
  \item \emph{Recall:} say the Latin for to seem, then the Latin for to serve, then say \emph{patior, patī} again
\end{itemize}

\begin{tcolorbox}[breakable,colback=teal!4,colframe=teal!35!black,title={Before you move on}]
What does patior, patī mean? (\textquotedblleft{}To suffer, to allow\textquotedblright{}.) Say it once more.
\end{tcolorbox}
